% !TeX root = main_long.tex
\section{VCF Core representation}
\label{sec:methods}

\label{sec:design}
\label{sec:profiles}

VCF Core is a semantic vocabulary with classes, properties, declared domains and ranges, and a shallow subclass hierarchy; it carries no OWL restrictions, cardinality, disjointness, or inverse axioms.
Value and cardinality constraints are expressed as SHACL shapes~\cite{shacl2017}.
The vocabulary was developed iteratively from the earlier VCF-RDFizer vocabulary~\cite{crum2025semantifying}, using AI-assisted drafting and author review to derive terms, field definitions, and constraints from VCF 4.1--4.5~\cite{htsspecs,vcf45}.
Terms and relationships can be browsed in the \href{https://ecrum19.github.io/vcf-core-vocabulary/}{VCF Core Vocabulary Explorer}~\cite{ocg2026}.

Listing~\ref{lst:vcf-example} shows a minimal VCF example illustrating the structure and typical field values contained in a VCF file.
The header declares a version, reference, contig, and three sample field definitions: genotype (\texttt{GT}), local alternate-allele indices (\texttt{LAA}), and local read depths (\texttt{LAD}).
The \texttt{FORMAT} column determines the identity of items within the colon-separated sample values, which two sample columns, \texttt{S1} and \texttt{S2}, then supply.
Unspecified VCF fields are shown as dots; trailing sample subfields may also be dropped entirely, as \texttt{S2} does on line 11.

\begin{lstlisting}[
  float=htbp,
  columns=fixed,
  basewidth={\dimexpr\linewidth/81\relax},
  numbers=left,
  numberstyle=\tiny,
  numbersep=8pt,
  caption={Four-record, two-sample VCF example (\texttt{example.vcf}), with columns aligned for display.},
  label={lst:vcf-example}
]
##fileformat=VCFv4.5
##reference=https://example.org/reference.fa
##contig=<ID=chr1,length=1000>
##FORMAT=<ID=GT,Number=1,Type=String,Description="Genotype">
##FORMAT=<ID=LAA,Number=.,Type=Integer,Description="Local ALT indices">
##FORMAT=<ID=LAD,Number=LR,Type=Integer,Description="Local read depths">
#CHROM POS ID REF ALT   QUAL FILTER INFO FORMAT     S1          S2
chr1   1   .  G   T     .    .      .    GT:LAA:LAD 0/1:1:18,22 0/0:1:21,0
chr1   2   .  G   A,C   .    .      .    GT:LAA:LAD 2/2:2:20,30 0/1:1:12,9
chr1   3   .  A   C,G,T .    .      .    GT:LAA:LAD 0/3:3:15,25 0/1:1:11,7
chr1   4   .  T   C     .    .      .    GT:LAA:LAD ./.:.:.     0/0
\end{lstlisting}

Figure~\ref{fig:model} shows how VCF Core models the fields of Listing~\ref{lst:vcf-example}, in three steps marked on the figure.
\stepmark{1} A \term{VCFFile} represents \texttt{example.vcf} and links to its \term{VCFHeader} (lines 1--6) and four \term{VCFRecord} resources (lines 8--11).
The header contains \term{HeaderLine} resources; a \term{FORMATHeaderLine} is also a \term{FieldDefinition}, supplying a field's identifier, type, and cardinality, as \texttt{LAD} does on line 6.
\stepmark{2} Each record describes a location and alleles; line 8 declares one \texttt{ALT} allele, \texttt{ALT=T}, modelled as a single \term{AltAllele}.
\stepmark{3} Its \term{VariantCall} groups field values and sample observations, and a \term{SampleCall} describes the values within the column named \texttt{S1} on line 7.

The first record is deliberately simple for \texttt{S1}: \texttt{GT=0/1} calls "G" and "T", and \texttt{LAA=1} selects the only \texttt{ALT} allele.
The declaration \texttt{Number=LR} means that \texttt{LAD=18,22} gives depth for the \texttt{REF} allele "G" followed by that local \texttt{ALT} allele "T"~\cite{vcf45}.
VCF Core represents "22" as a \term{FieldValueItem} linked through \term{forAllele} to "T", and through its \term{FormatFieldValue} to the declaration and sample observation.


\begin{figure}[htbp]
\centering
\begin{tikzpicture}[x=1mm,y=1mm,
  box/.style={draw,rounded corners=1pt,align=center,font=\small,inner sep=4pt},
  expanded/.style={box,fill=blue!10},
  condensed/.style={box,fill=orange!30},
  edge/.style={-{Stealth[length=4pt]},thin},
  lab/.style={font=\scriptsize\ttfamily,fill=white,inner sep=1pt,align=center}]
\node[box] (file) at (17,0) {VCFFile\\\scriptsize\texttt{example.vcf}};
\node[box] (header) at (72,0) {VCFHeader\\\scriptsize lines 1--6};
\node[box] (line) at (128,0) {FORMATHeaderLine\\\scriptsize\texttt{LAD, Number=LR}\\\scriptsize line 6};
\node[box] (record) at (17,-18) {VCFRecord\\\scriptsize line 8};
\node[box] (allele) at (72,-18) {AltAllele\\\scriptsize\texttt{alleleValue "T"}\\\scriptsize line 8, \texttt{ALT}[1]};
\node[expanded] (item) at (128,-18) {FieldValueItem\\\scriptsize\texttt{itemValue 22}\\\scriptsize line 8, \texttt{LAD}[2]};
\node[box] (call) at (17,-36) {VariantCall};
\node[expanded] (sample) at (72,-36) {SampleCall\\\scriptsize\texttt{sampleId "S1"}\\\scriptsize line 7};
\node[expanded] (value) at (128,-36) {FormatFieldValue\\\scriptsize\texttt{fieldValue "18,22"}\\\scriptsize line 8, \texttt{LAD}};
\node[condensed] (matrix) at (72,-54) {CohortCallMatrix};
\node[condensed] (vector) at (128,-54) {FormatValueVector\\\scriptsize\texttt{encodedValues "18,22"}\\\scriptsize line 8, \texttt{LAD}};
\node[anchor=east] at ([xshift=-2mm]file.west) {\stepmark{1}};
\node[anchor=east] at ([xshift=-2mm]record.west) {\stepmark{2}};
\node[anchor=east] at ([xshift=-2mm]call.west) {\stepmark{3}};
\draw[edge] (file) -- node[lab,above] {hasHeader} (header);
\draw[edge] (header) -- node[lab,above] {hasHeaderLine} (line);
\draw[edge] (file) -- node[lab,left] {hasRecord} (record);
\draw[edge] (record) -- node[lab,left] {hasCall} (call);
\draw[edge] (record) -- node[lab,above] {hasAltAllele} (allele);
\draw[edge] (item) -- node[lab,above] {forAllele} (allele);
\draw[edge] (value) -- node[lab,right] {hasValueItem} (item);
\draw[edge] (call) -- node[lab,above] {hasSampleCall} (sample);
\draw[edge] (sample) -- node[lab,above] {hasFormatValue} (value);
\draw[edge] (call.south) |- node[lab,above,pos=0.68] {hasCallMatrix} (matrix);
\draw[edge] (matrix) -- node[lab,above] {hasFormat\\ValueVector} (vector);
\draw[edge,dashed] (vector.east) -- (158,-54) |- (line.east);
\draw[dashed,thin] (value.east) -- (158,-36);
\node[lab,rotate=90] at (158,-23) {declaredBy};
\end{tikzpicture}
\caption{
  Selected resources for the first record (line 8) of Listing~\ref{lst:vcf-example}, walked through in steps \stepmark{1}--\stepmark{3}. 
  Purple nodes are \textit{expanded}-profile resources, orange \textit{condensed}. 
  Depth 22 is explicitly associated with allele T in the \textit{expanded} profile. 
  \texttt{LAA} has no node of its own because its resolution is already applied in \term{forAllele}. 
  The two dashed links show that both profiles resolve to the same \term{FieldDefinition}. 
  Records on lines 9--11, column \texttt{S2}, other fields, alleles, and header lines are omitted.
  }
\label{fig:model}
\end{figure}

VCF Core provides two \emph{sample profiles} for representing sample-level observations.
In the \emph{expanded} profile, each sample value is represented as its own RDF resource. 
In the \emph{condensed} profile, values are stored as sample-ordered vectors, avoiding per-sample field resources but requiring decoding to access an individual value.

The difference is primarily one of scale for multi-sample VCF files (as used in cohort studies such as~\cite{1000genomes2015global}). 
\[
\underbrace{O\!\left(N_{\mathrm{variants}}N_{\mathrm{samples}}N_{\mathrm{FORMAT}}\right)}_{\text{(E) Expanded}}
\qquad
\underbrace{O\!\left(N_{\mathrm{variants}}N_{\mathrm{FORMAT}}\right)}_{\text{(C) Condensed}} .
\]

For illustration, a VCF with \(1{,}000\) variants (rows) across \(2{,}500\) samples (sample columns) using the \texttt{FORMAT} column \texttt{GT:LAA:LAD} from above, would produce approximately \(1.1\times10^{8}\) sample-related statements in expanded profile (e), and approximately \(3.05\times10^{4}\) in condensed profile (c).
For datasets containing millions of variants and thousands of samples, reducing the number of triples by several orders of magnitude using the condensed profile can substantially decrease graph size and make storage and usage more practical.
Importantly, both profiles preserve the same underlying values, the trade-off is direct queryability. 
Expanded graphs support SPARQL access to individual sample values and their relationships, while condensed graphs require a decoding step before those values can be queried individually.
In the condensed profile, values are stored inside vector literals and recovered using the corresponding FORMAT definition and sample ordering, see \href{https://github.com/ecrum19/VCF-RDFizer/blob/main/docs/sample-representation-guide.md}{github docs} for more detail\footnote{\url{https://github.com/ecrum19/VCF-RDFizer/blob/main/docs/sample-representation-guide.md}}.
Both profiles retain the same file, header, and record context, including version-specific registries and validation overlays selected from \term{fileFormat}~\cite{vcfc210}.


\subsection{Assessment design}
Each research question adapts an established practice.
RQ1 states each specification rule as a question with a checkable answer~\cite{gruninger1995methodology} and tests it after SAMOD~\cite{samod2016}, while its curated inventory evaluates the vocabulary against an external frame of reference~\cite{gomezperez2001evaluation}.
RQ2 adapts conformance testing, executing the same specification-derived tests against independent implementations, as done in other fields~\cite{bishop2005conformance} and for tools reading one genomics format~\cite{niu2022interop}.
RQ3 and RQ4 are task-based and comparative evaluations, judging the representation by whether it supports an external task and against related models through their published artifacts~\cite{sabou2012ontologyevaluation}.

\paragraph{Evaluated artifacts and reproduction.}
\label{sec:reproduction}
All results use VCF Core v2.1.3~\cite{vcfc210} and VCF-RDFizer v3.0.3, with the latter's container pinned by digest.
The \href{https://github.com/ecrum19/SWAT4HCLS_2027}{experimental repository}\footnote{\url{https://github.com/ecrum19/SWAT4HCLS_2027}} provides the referenced files, requirements, fixtures, queries, expected answers, per-RQ procedures, and pinned versions and digests in \path{PINS.md}.
Runs record specification digests, commands, graphs, results, and environments, including rdflib~7.6.0 for query execution.
